\section{A bounded region of changing tree attachments}
\label{sec:op3-bounded-attachments}

\paragraph{Status and model.}
The following are \statusproved{} proof drafts, awaiting independent review.
The graph is a simple cycle with disjoint finite trees attached to its
centers through one edge each. The seed is on the cycle. Neither the cycle
ordering nor any attachment size is supplied. Use the degree obstacle
problem with $\bm M=\bm D-\gamma\bm A$,
$\bm b=\bm e_v-\lambda\bm d$, $0<\gamma<1$, and $\lambda>0$.
Let $q\geq1$ be the maximum size of an attachment incident to a cycle
center that is positive in the final answer, taking $q=1$ if there are no
such attachments. This is an analysis parameter, not advice to the solver.
The exact-real word model charges all bounded probes, even into inactive
vertices, all curve construction, stored-state work, event searches and
output. No coefficient-bit or floating-point claim is made.

\subsection{When a tree attachment has finished changing}

\begin{lemma}[A discounted-hitting saturation bound]
\label{lem:op3-attachment-saturation}
Let a tree $T$ of at most $q$ vertices be attached through one edge to a
center whose potential is fixed at $x\geq0$. All loads in $T$ are
$-\lambda d_i$, using actual graph degrees. Its conditional obstacle
solution has every coordinate strictly positive, and its response has
reached the final affine piece, whenever
\begin{equation}
 x\geq X_q:=\frac{\lambda}{1-\gamma}
                    \bigl(\gamma^{-2q^2}-1\bigr).
 \label{eq:op3-attachment-saturation}
\end{equation}
The algorithm need not evaluate this threshold.
\end{lemma}
\begin{proof}
Run simple random walk until it first hits the center, with hitting time
$\tau$, and put $p_i=\mathbb E_i[\gamma^\tau]$. Harmonic conditioning on
the first step gives $p_i=\gamma\sum_{j\sim i}p_j/d_i$, with center value
one. Direct substitution therefore shows that the unconstrained solution is
\begin{equation}
 u_i=p_i\left(x+\frac{\lambda}{1-\gamma}\right)
                -\frac{\lambda}{1-\gamma}.
 \label{eq:op3-discounted-attachment}
\end{equation}

Root the attachment at the center. The expected time from a child $j$ to
its parent is $2|T_j|-1$. Indeed, first-step decomposition gives
$H_j=d_j+\sum_{k\text{ child of }j}H_k$; induction over the descendant
tree gives the displayed expression. The walk from any $i$ to the center
crosses the successive parent edges in order. The strong Markov property
then gives
$\mathbb E_i\tau\leq q(2q-1)<2q^2$.
The function $t\mapsto\gamma^t$ is convex, so
$p_i\geq\gamma^{\mathbb E_i\tau}>\gamma^{2q^2}$.
Under \cref{eq:op3-attachment-saturation},
\cref{eq:op3-discounted-attachment} is strictly positive for every $i$.
It is therefore the obstacle optimum. It remains positive at larger $x$,
so the response thereafter is this single affine formula.
\end{proof}

\begin{lemma}[Potential growth behind the active end]
\label{lem:op3-cycle-interior-growth}
Before cycle closure, let $y_0\geq0$ be the potential at an active end,
$y_{-1}=0$ the next inactive cycle coordinate, and $y_k$ the consecutive
active-center potentials toward the seed. Up to and including the seed,
\begin{equation}
 y_k\geq\frac{\lambda}{1-\gamma}(\gamma^{-k}-1).
 \label{eq:op3-cycle-interior-growth}
\end{equation}
Consequently all attachments at a center at least $2q^2$ edges behind an
active end have finished changing.
\end{lemma}
\begin{proof}
A negative-load attachment has no coordinate greater than its center:
a larger positive maximum would contradict its own obstacle equality and
$\gamma<1$. At a nonseed cycle center of degree $d_i=2+k_i$, the sum of
the $k_i$ attached-root potentials is at most $k_i y_i$. Its active
equation therefore implies
\[
 y_{i+1}+y_{i-1}
 =\frac{2+k_i}{\gamma}(y_i+\lambda)
     -\sum_{j\text{ attached at }i}u_j
 \geq\frac2\gamma(y_i+\lambda).
\]
Put $w_i=y_i+\lambda/(1-\gamma)$. Then
$w_{i+1}\geq2w_i/\gamma-w_{i-1}$, and
$w_0\geq w_{-1}=\lambda/(1-\gamma)$.
The first step gives $w_1\geq w_0/\gamma$. If
$w_i\geq w_{i-1}/\gamma$, the recurrence gives
$w_{i+1}\geq(2/\gamma-\gamma)w_i\geq w_i/\gamma$.
Induction proves \cref{eq:op3-cycle-interior-growth}. This uses only
nonseed equations even when the last derived value is the seed value.
Apply \cref{lem:op3-attachment-saturation} at distance $2q^2$.
\end{proof}

The length-two counterexample in
\cref{prop:op3-cycle-unsettled-attachment} is consistent with these bounds.
An attachment can still change after one advance; it cannot remain
unsettled arbitrarily far behind the active end when its size is bounded.

\subsection{Finding ports and building only bounded components}

At an admitted cycle center $c$, read its adjacency row and try exploration
limits $L=1,2,4,\ldots$. For each neighbor $j$, explore its component in
$G-c$ until it is exhausted or one of the following certificates occurs:
more than $L$ vertices are found, a vertex has degree greater than $L$, or
another neighbor of $c$ is reached. In the degree test, stop before reading
that vertex's adjacency row. An exhausted component is an attached tree of
size at most $L$. The other cases provisionally mark a port. Reaching
another neighbor marks both of the corresponding ports.

The actual two cycle ports are always marked: their common component in
$G-c$ either exceeds the limit or connects the two neighbors. An attached
tree can be marked only when it exceeds $L$. Thus exactly two marked
ports certify that all attached components have been fully discovered.
Otherwise double $L$ and retry, caching graph replies. If the largest
attachment at $c$ has size $q_c\geq1$, this terminates with $L<2q_c$,
or with $L=1$ when there is no attachment. Each trial explores at most
$L$ rows of degree at most $L$ per incident neighbor, plus a constant
number of stopping replies. Its charged work is $O(L^2d_c)$; all trials
sum to $O(q_c^2d_c)$. No ambient graph scan or structural advice is used.

Build each discovered attachment's response once using
\cref{thm:op3-persistent-tree}, retaining its child curves for terminal
reconstruction. It is permissible to build already exhausted small
components again during failed classification trials, provided this work
is charged; the implementation does so. Each attachment response is
convex and has at most its size in knots. For a center, sort the union of
its attachment knots in the center's own potential and maintain one cursor.
If its current summed response is $C_i x+E_i$, the effective center
diagonal and load are $d_i-C_i$ and $-\lambda d_i+E_i$; the seed also has
its source load one. Crossing a hinge $c(x-\tau)_+$ subtracts $c$ from
that diagonal and $c\tau$ from that load. Equal knots are consumed together.

\subsection{Eagerly store settled prefixes and retain the changing region}

Run the exact source-potential continuation of
\cref{sec:op3-local-cycle}. On each arm retain a consecutive window of
centers adjacent to the active end. After any event, repeatedly eliminate
the first center of the window if its attachment cursor is exhausted,
always leaving at least the last center. Eliminate from the seed outward,
store one fixed reverse record, and accumulate its source-derivative
contribution as in \cref{eq:op3-cycle-append,eq:op3-cycle-root-derivative}.
The rule depends on the exhausted cursor, not on knowing $q$.

If a window has more than one center, its first center has an unsettled
attachment. By \cref{lem:op3-cycle-interior-growth}, its distance to the
active end is less than $2q^2$. Therefore each retained arm has at most
$2q^2$ centers. The attached components have already been compressed to
their current scalar affine pieces. Their remaining cycle system is
tridiagonal, with affine loads in the source potential $t$. Two forward
and reverse passes compute all retained response coefficients in $O(q^2)$
work. The exact root derivative follows from their Schur contributions.

Inspect only the next attachment knot at each retained center, the seed's
next knot, and the one or two cycle candidates. Every knot time is the
zero of a known affine equality. The earliest event or earlier root of
the reduced derivative is therefore found in $O(q^2)$ work, including
the no-event certificate. A fully settled eliminated center has no future
attachment event. No old prefix coordinate is queried or rewritten.

If the two candidates coincide, use their combined load for the closing
gate. After its admission, the two windows and the last center form one
tridiagonal path, ordered from one seed-side prefix to the other, of size
at most $4q^2+1$. Retain this path for the rest of the continuation.
It contains every remaining nonseed attachment event. The root's own
knots are still checked directly. Compute the two affine right-hand sides
and the first remaining event by linear passes in this bounded path.
Finish with one reverse pass through all stored prefixes and one traversal
of the retained attachment curves.

\begin{theorem}[Local solve with bounded tree attachments; proof draft]
\label{thm:op3-bounded-attachments}
For a cycle-seeded graph in the stated family, the described algorithm
returns the exact obstacle optimum with
\begin{equation}
 \widetilde O\bigl(q^3(1+\vol(S))\bigr)
 \quad\text{word work and}\quad
 \widetilde O\bigl(q^2(1+\vol(S))\bigr)
 \quad\text{working storage},
 \label{eq:op3-bounded-attachment-work}
\end{equation}
where $S$ is its positive support. The quantity $q$ and the graph structure
need not be supplied. Admitted center rows and all bounded probes are
charged. Taking $\lambda=\epsappr$ gives an ACL approximation with work
$\widetilde O(q^3/\epsappr)$ and no polynomial dependence on $1/\alpha$.
For every fixed attachment-size bound this attains the OP3 scale on the
promised family. It is not an arbitrary-graph or an arbitrary-seed theorem.
\end{theorem}
\begin{proof}
Port classification is exact under the graph promise. At each admitted
center the supplied bounded components therefore give its entire exact
conditional response. The two-arm description enumerates all cycle gates
before closure, including combined load at the last inactive center.
All other possible events are attachment knots. Eliminated centers have
exhausted their lists; the retained windows and the seed contain the rest.
The source derivative, response equations and knot times are exact Schur
computations on these current pieces. The positive Schur derivative and
continuity justify advancing precisely when its zero lies after the next
event. Thus every admitted center is positive at the final root; equality
causes stopping before a new row is admitted. Final reconstruction solves
all active equations and all other obstacle inequalities.

Summing center degrees is bounded by $\vol(S)$. Classification, including
failed bounded probes and reaccesses, costs $O(q^2\vol(S))$. The number
of attachment vertices built is at most $q$ times the number of attached
components, hence $O(q\vol(S))$. Repeated bounded builds and sorting cost
$\widetilde O(q\vol(S))$. The number of consumed knots and center
admissions is $O(q\vol(S))$. Each costs $O(q^2)$ for retained factors,
responses and event searches, including unsuccessful searches. Each
committed record is created once and used once at the end; even array
shifts of retained windows cost only $O(q^2)$ per committed record.
Terminal attached-curve evaluations cost
$\widetilde O(q\vol(S))$ and positive output writes cost $O(\vol(S))$.
These terms give the work bound. Cached probe incidences, persistent
bounded-component curves, current knot lists, center records and temporary
tridiagonal arrays fit the stated working storage bound; discarded
construction and solve allocations are already charged to work.

Finally the full-graph obstacle equations give
$\lambda\vol(S)\leq1$, and the inverse-positive residual argument of
\cref{cor:op3-local-cycle-acl} gives the ACL error
$\bm0\leq\bm\pi-\bar\alpha\bm D\bm u\leq\lambda\bm d$.
\end{proof}

\paragraph{Measured evidence and remaining question.}
The exact bounded-attachment implementation passes 1,909
comparisons against a separate exact obstacle solver, including branching
attachments and the preceding length-two witness. Eight larger instances
pass full-graph KKT checks, performed only by the audit. There are also
195 exact discounted-hitting checks on rooted trees through size six.
For sizes two and three the reported larger runs retain at most five cycle
centers at once; the theorem uses the conservative bound $4q^2+1$.
The implementation counts failed classification probes, repeated curve
construction and every retained-window factor row. These finite checks do
not replace independent proof review. The remaining OP3 question is to
remove the growing-attachment factor or replace it by a paid approximate
response certificate, and then to handle more than one cyclic core.
